\subsection{Henselization: source review and exact comparison maps}\label{algebra-proof-060-henselization-source-review-and-exact-comparison-maps}

Primary source: Stacks Project authors, \texttt{algebra.tex}, original authority commit \texttt{a04446e57ec1fbc252a871afcec7752fb2807b14}, lines 42959--43572, file SHA256 \texttt{FA8BB92E58A4F78A2BD01B3B6A4A87DE0A0D279F5DD90641B574DD5FBFFFA4F3}. This is a separate editorial validation of the original statements and their omitted details. It does not replace the source exposition or claim a new theorem. The source reading includes the next opening, 43573--43593; that opening remains unadjudicated.

The user's latest direction keeps the Algebra review first. Optional automorphism, invariant-ring and more general residue-extension investigations are deferred, along with additional paper mining and the later synthesis. Their consideration creates no proved-consequence group here. The eight minimal edits below account for fifteen received reports in nine groups.

The previously proved inputs used below are the finite-presentation descent and étale permanence statements, the henselian finite-algebra decomposition and equivalence of finite étale categories with residue-field finite étale categories, and the exact ind-étale comparison and henselian-target lifting statements. Their full arguments are in \texttt{ALGEBRA\_ETALE\_NOTES\_20260928.md}, \texttt{ALGEBRA\_HENSELIAN\_NOTES\_20260929.md} and \texttt{ALGEBRA\_INDETALE\_NOTES\_20260929.md}; the evidence receipt binds their bytes. These are mathematical inputs, not assumptions added to the source's hypotheses.

\subsubsection{Pointed systems and the original coequalizer}\label{algebra-proof-060-pointed-systems-and-the-original-coequalizer}

Let \((R,\mathfrak m,\kappa)\) be the original local ring. An ordinary object is precisely the source pair \((S,\mathfrak q)\), with \(S/R\) étale and its structural residue map \(\kappa\to\kappa(\mathfrak q)\) an isomorphism. Every arrow is an \(R\)-map contracting the selected target prime to the selected source prime. In the strict system a fixed inclusion \(\kappa\subset\kappa^{sep}\) is retained, and the original triple carries an embedding \(\alpha:\kappa(\mathfrak q)\to\kappa^{sep}\). An arrow must preserve this embedding. This condition is essential data even when the underlying arrow is also an arrow of pairs.

The systems are nonempty, containing \(R\) with its selected prime and residue map. For two ordinary objects the source ring is exactly \(S''=S\otimes_R S'\), with \(\mathfrak q''=\mathfrak qS''+\mathfrak q'S''\). Since \(\mathfrak m\) is maximal and \(S\) is étale, its closed fibre is a finite product of finite separable fields. In the ordinary case the selected factor is \(\kappa\); thus \(S/\mathfrak q=\kappa=S'/\mathfrak q'\) and \(S''/\mathfrak q''=\kappa\). Both structure maps contract \(\mathfrak q''\) as required.

For two strict objects use the same original tensor product, and the residue map \[
 S\otimes_R S'\longrightarrow\kappa^{sep},\qquad
 x\otimes y\longmapsto\alpha(\overline x)\alpha'(\overline y).
\] Its kernel is the selected prime. Its residue field is the compositum of the two embedded finite separable fields: the image contains both fields and is their finite-dimensional tensor image, a domain finite-dimensional over \(\kappa\), hence a field. This proves the residue assertion and both required contractions. Étaleness is preserved by this tensor product.

For parallel arrows \(\varphi,\psi:S\to S'\), keep the source's full ring, rather than replacing it in the construction: \[
 U=(S'\otimes_{\varphi,S,\psi}S')\otimes_{S'\otimes_R S'}S'.
\] The map from \(S'\otimes_R S'\) to the first factor is the canonical quotient map, and the map to the last factor is multiplication. Let \(J\subset S'\) be the ideal generated by every \(\varphi(s)-\psi(s)\), for \(s\in S\). The exact comparison maps are \[
 U\longrightarrow S'/J,\quad (a\otimes b)\otimes c\longmapsto[abc],
 \qquad
 S'/J\longrightarrow U,\quad[c]\longmapsto(1\otimes1)\otimes c.
\] The first map respects the \(S\)-balancing because the difference of the two images is \([ab(\varphi(s)-\psi(s))c]=0\). It respects the \(S'\otimes_R S'\)-balancing because multiplication of the two factors has the same image \([abc]\). In \(U\), balancing gives \((a\otimes b)\otimes c=(1\otimes1)\otimes abc\), and the original relation \(\varphi(s)\otimes1=1\otimes\psi(s)\) therefore kills \(J\) under the second map. The two displayed maps compose to the identity on every tensor generator and every quotient class. In particular the original map \(S'\to U\) equalizes \(\varphi\) and \(\psi\).

All étaleness claims here follow on the original factors. The arrows \(S\to S'\) are maps between étale \(R\)-algebras, hence étale. The first tensor factor is étale over \(S'\) by base change and over \(R\) by composition. Multiplication \(S'\otimes_R S'\to S'\) is also a map between étale \(R\)-algebras, hence étale; its base change to the first factor makes \(U\) étale over that factor and over \(R\).

Write \(\rho':S'\to\kappa\), or \(S'\to\kappa^{sep}\), for the selected residue map. The pair/triple conditions give \(\rho'\varphi=\rho'\psi\). Thus the exact map \[
 U\longrightarrow\kappa\text{ or }\kappa^{sep},\qquad
 (a\otimes b)\otimes c\longmapsto\rho'(a)\rho'(b)\rho'(c)
\] is defined. Its kernel contracts to \(\mathfrak q'\) in the last factor. In the ordinary local construction its kernel is exactly the source ideal \[
 (\mathfrak q'\otimes S'+S'\otimes\mathfrak q')\otimes S'
 +(S'\otimes_{\varphi,S,\psi}S')\otimes\mathfrak q'.
\] Indeed the comparison with \(S'/J\) sends this ideal to \(\mathfrak q'/J\), since \(J\subset\mathfrak q'\), and the quotient is \(S'/\mathfrak q'=\kappa\). The strict construction has the same contraction and the inherited embedding of \(\kappa(\mathfrak q')\). This proves filteredness, including the parallel-arrow condition, for the original systems.

For a general base prime \(\mathfrak p\subset R\), use the source fibre rings \(F=S\otimes_R\kappa(\mathfrak p)\) and \(F'=S'\otimes_R\kappa(\mathfrak p)\). They are finite products of finite separable fields. Their selected factor maps define precisely the same product and coequalizer residue maps. The original source formula is \[
 F''=(F'\otimes_{\varphi,F,\psi}F')
       \otimes_{F'\otimes_{\kappa(\mathfrak p)}F'}F'.
\] The selected map \(F'\to\kappa(\mathfrak p)\) belongs to \(\mathfrak q'\), not to an undefined \(\mathfrak p'\). Its extension sends \((a\otimes b)\otimes c\) to the product of the three selected images. This proves the correction at 43284 on the exact original ring. In the strict case its target is \(\kappa^{sep}\), and the triple compatibility gives the same balancing identity. The selected prime of the original ring is the inverse image of the fibre prime. One must not assume that the unlocalized map \(S\to\kappa(\mathfrak q)\) is surjective for a general base prime; its image has fraction field \(\kappa(\mathfrak q)\).

\subsubsection{Local colimits and the source henselian property}\label{algebra-proof-060-local-colimits-and-the-source-henselian-property}

For any selected \(\mathfrak q\) and any \(f\in S\setminus\mathfrak q\), the original system includes \(S_f\) with its selected prime and, in the strict case, the unchanged residue embedding. Hence every such \(f\) becomes invertible in the colimit \(H\). The canonical maps give inverse comparisons \[
 H=\mathop{\rm colim}S\longrightarrow\mathop{\rm colim}S_{\mathfrak q},
 \qquad s/f\longmapsto [s][f]^{-1}\in H.
\] The latter formula respects every localization relation and transition map; the composites fix \([s]\) and every \(s/f\). This proves the two colimit formulas with the original objects retained. In the localized system every transition map is local. The colimit of their residue fields is \(\kappa\) in the ordinary case and \(\kappa^{sep}\) in the strict case. In the latter assertion each element of \(\kappa^{sep}\) lies in a finite separable extension of \(\kappa\), and that extension occurs as a selected étale residue field by the previously proved prescribed-residue construction. Common tensor objects put every finite list in a single stage.

At each localized stage, étaleness gives \(\mathfrak pR_{\mathfrak p}S_{\mathfrak q}=\mathfrak qS_{\mathfrak q}\). The kernel of the colimit residue map is exactly \(\mathfrak pH\): an element with zero residue has zero residue already at its representing stage, since the residue transition maps are field embeddings, and then belongs to the displayed stage ideal. Conversely that ideal maps to zero. An element outside this kernel is represented by a unit in a localized stage. Thus \(H\) is local, its maximal ideal is precisely \(\mathfrak pH\), and its specified residue field is the indicated field colimit. This also proves that the zero ring cannot occur as a selected stage or final local ring.

For the original local-base argument at 43040--43056, all primes of the closed fibre are maximal, not only the selected prime. Distinct selected and unselected primes are therefore incomparable. Choose \(g_i\in\mathfrak q_i\setminus\mathfrak q\) and keep the full product \(g=\prod_{i=1}^r g_i\); for an empty list take \(g=1\). The localization at \(g\) removes every other closed-fibre factor. That fibre is now its single selected field, so \(\mathfrak mS_g=\mathfrak qS_g\), not merely an equality of radicals. This justifies the source's stronger ideal equality. The source inverse of \([f]\) is exactly the class of \(1/f\) in \(S_f\).

For the henselian property, let the original monic polynomial \(P\in H[T]\) and its specified simple residue root \(a_0\) be given. Filteredness places all coefficients at one stage, with a monic lift \(Q\). In the ordinary local construction \(a_0\in\kappa=S/\mathfrak q\), so choose its original lift \(a'\in S\). In the strict construction first enlarge to a common pointed stage containing the finite separable residue extension generated by \(a_0\), together with the coefficient stage. Since this stage lies over the maximal ideal of the local base, its selected prime is maximal and its quotient is its residue field; therefore \(a_0\) also has a lift \(a'\) there. This is the additional residue step implicit in the source's reference to the same proof.

Keep exactly \(S'=S[T]/(Q)\) and \(\mathfrak q'=\mathfrak qS'+(T-a')S'\). The quotient is the selected residue field. The derivative \(Q'(T)\) is nonzero there, since \(a_0\) is a simple root. Localizing at its class gives an étale \(S\)-algebra and preserves \(\mathfrak q'\); equivalently use the source's chosen \(g\notin\mathfrak q'\) in the étale locus. The class of \(T\) in that object maps to an element \(a\in H\) satisfying the exact equations \(P(a)=0\) and \(\overline a=a_0\). This proves the henselian property by the simple-root criterion already established in the preceding section. In the strict case the residue field is separably closed, so \(H\) is strictly henselian.

For an arbitrary \(R,\mathfrak p\), the source's descent from \(R_{\mathfrak p}\) uses the proved finite-presentation descent for étale algebras: each étale algebra over \(R_{\mathfrak p}\) descends to one over \(R_f\) for some \(f\notin\mathfrak p\). The composite from \(R\) is étale; contraction of the selected prime gives an object over \((R,\mathfrak p)\) with unchanged residue field. Conversely the base localization of every such object is an object over the local base. Every finite set of algebra generators, arrow images, and relations of two arrows descends after a further denominator \(g\notin\mathfrak p\), by finite presentation. Thus not only the objects but the equalities used in the colimit descend. The inverse comparison sends a representative at a descended object to its original colimit class; independence of the descent follows at a common further localization. This proves both bijectivity and ring compatibility of the source's canonical ordinary and strict comparisons. No injectivity of transition maps is assumed in these comparisons.

\subsubsection{Residue choices and functoriality}\label{algebra-proof-060-residue-choices-and-functoriality}

For a henselian local target \(B\), an étale \(R\)-algebra \(A\), a selected prime over the contracted maximal ideal, and a specified compatible embedding of its residue field into \(\kappa(B)\), the previously proved henselian lifting lemma gives one and only one map \(A\to B\) inducing these data. On a filtered system the maps are compatible: for each transition arrow the two candidate composites have identical base and residue maps, so stagewise uniqueness identifies them. The resulting colimit map is unique and local, because its reduction is the specified field embedding and the source maximal ideal is the kernel of its specified reduction. This proves all the functoriality assertions in 43186--43224 and 43337--43397 on the original systems.

In particular, the map in 43353 has target \(S^{sh}\). Its correctly typed contraction is \[
 f^{-1}(\mathfrak m_{S^{sh}})=\mathfrak q.
\] The source's \(\mathfrak m_{S^h}\) at 43360 belongs to a different ring, so the proposed replacement is a mathematical target correction. The residue condition still uses precisely the original chosen map \(\phi:\kappa(\mathfrak q)\to\kappa^{sep}\); no map has been discarded.

The unique-isomorphism assertions at 43143--43148 and 43157--43162 are understood in their specified residue-field data. Given two candidate rings with those data, the lifting lemma constructs maps in both directions, and uniqueness makes both composites identities. For the strict construction the given identification with \(\kappa^{sep}\) is part of that comparison, just as the arrows in the source explicitly preserve residue maps. This is a contextual reading of the source, not a new assertion that arbitrary residue automorphisms must be trivial. Investigation of the larger automorphism and fixed-ring structures remains deferred.

For the ordinary base-change statement, put \(H=(R_{\mathfrak p})^h\), \(K=(S_{\mathfrak q})^h\) and retain \(T=H\otimes_R S\). Its original comparison map is \[
 T\longrightarrow K,\qquad h\otimes s\longmapsto f(h)s.
\] The prime \(\mathfrak t\) is the kernel of its composite to \(\kappa(\mathfrak q)\), given by \(h\otimes s\mapsto\overline h\,\overline s\) using the structural residue extension. It is the unique prime contracting to \(\mathfrak m_H\) and \(\mathfrak q\): primes with those contractions correspond to primes of \[
 \kappa(\mathfrak p)\otimes_{\kappa(\mathfrak p)}\kappa(\mathfrak q)
 =\kappa(\mathfrak q),
\] which has just its zero prime. This correspondence is obtained by quotienting the two maximal/prime ideals and inverting all nonzero elements in both residue domains. The image of the unlocalized ring \(T\) need not be all of this field, but its fraction field is \(\kappa(\mathfrak q)\), because the image already contains \(S/\mathfrak q\). Hence this is exactly \(\kappa(\mathfrak t)\).

Both \(T\) and \(K\) are ind-étale over \(S\), so the original comparison makes \(K\) ind-étale over \(T\), by the earlier maps-between-ind-étale-algebras result. Every element outside \(\mathfrak t\) maps to a unit of \(K\); base change along \(T\to T_{\mathfrak t}\) consequently retains \(K\) itself. Thus \(K/T_{\mathfrak t}\) is ind-étale and local with equal residue field, and \(K\) is henselian. The henselian uniqueness lemma identifies it with the henselization of \(T_{\mathfrak t}\). The two inverse maps are the unique local maps inducing the identity on the common residue field; their composites are identities by that same uniqueness. This proves the source's identification, rather than only an abstract existence assertion.

For strict base change retain \(H^{sh}=(R_{\mathfrak p})^{sh}\), \(K^{sh}=(S_{\mathfrak q})^{sh}\), \(T^{sh}=H^{sh}\otimes_R S\), and the given \(\phi:\kappa_1^{sep}\to\kappa_2^{sep}\). The selected prime is precisely \[
 \mathfrak t_\phi=\ker\bigl(T^{sh}\longrightarrow\kappa_2^{sep},
            \ h\otimes s\longmapsto\phi(\overline h)\overline s\bigr).
\] Its residue field is the compositum \(L=\kappa(\mathfrak q)\,\phi(\kappa_1^{sep})\) inside \(\kappa_2^{sep}\). Indeed the fraction field of the image contains both indicated fields and is generated by them. The extension \(L/\kappa(\mathfrak q)\) is algebraic separable: each finite list of elements from \(\phi(\kappa_1^{sep})\) lies in a finite separable extension of \(\kappa(\mathfrak p)\), and after base change the selected field factor is finite separable over \(\kappa(\mathfrak q)\). Every element of \(\kappa_2^{sep}\) is algebraic separable over \(L\), because its separable minimal polynomial over \(\kappa(\mathfrak q)\) remains square-free and its minimal polynomial over \(L\) divides it. Since \(\kappa_2^{sep}\) is separably closed, it is therefore a separable closure of \(L\). The map to \(K^{sh}\) is ind-étale, and remains so after localization at \(\mathfrak t_\phi\), by the same exact base-change argument. The target is strictly henselian with the specified residue closure. Uniqueness with that residue embedding gives the source's strict-henselization identification and the inverse map. The source correctly says a prime, not a unique prime independent of \(\phi\).

\subsubsection{The finite étale tower and the original tensor formula}\label{algebra-proof-060-the-finite-uxe9tale-tower-and-the-original-tensor-formula}

Let \(A=(R_{\mathfrak p})^h\). For each finite separable subextension \(\kappa'\subset\kappa^{sep}\), the finite étale equivalence for the henselian ring \(A\) gives the local finite étale algebra \(A(\kappa')\), together with its specified residue identification. Every inclusion \(\kappa'\subset\kappa''\) lifts to a unique map \(A(\kappa')\to A(\kappa'')\). Uniqueness gives the identity and composition laws exactly. This map is finite because the target is finite over \(A\), and étale because it is a map between étale \(A\)-algebras. It is local: the maximal ideals are \(\mathfrak m_A A(\kappa')\) and \(\mathfrak m_A A(\kappa'')\), and the residue map is an injection. A flat local map is faithfully flat, hence injective. Consequently the source's union denotes an actual system of injective maps, not an unjustified injectivity assumption about the earlier arbitrary stages.

For completeness the faithful-flat injectivity used here follows on the original module map. Flatness sends \(I=\ker(A_1\to A_2)\subset A_1\) to an injection \(I\otimes_{A_1}A_2\to A_2\), whose image is zero; faithful flatness gives \(I=0\). Flatness is faithful for a local map because every proper ideal \(J\subset A_1\) is contained in its maximal ideal, and \(JA_2\) is contained in the target maximal ideal, so \(A_2/JA_2\ne0\). The established flatness criterion by quotients then gives faithful flatness.

Each tower stage is henselian, by the finite-over-henselian theorem. Their filtered colimit with local maps is local with residue \(\kappa^{sep}\) and is henselian. To verify the last assertion directly, put the coefficients of a monic polynomial and its selected residue root in one stage. The residue embeddings ensure that the derivative is already nonzero at that stage when it is nonzero in the limit. The stage's henselian property gives a root there with that residue, and its image is the required root in the original colimit. The colimit is ind-étale over \(A\), and by composition over \(R\). The specified residue comparison and uniqueness identify the tower colimit with the original \(R^{sh}\).

This also supplies all the omitted details at 43488--43505 without assuming that inclusion of an indexing subsystem implies injectivity on its colimit. Let \(C\) be the original triple colimit in that passage. The preceding colimit proof makes \(C\) local, with maximal ideal \(\mathfrak pC\) and residue \(\kappa^{sep}\). The ordinary subsystem gives a map \(A\to C\) compatible with \(\kappa\subset\kappa^{sep}\), hence local. Both \(A\) and \(C\) are ind-étale over \(R\); therefore \(C\) is ind-étale over \(A\). This map is flat and local, hence faithfully flat and injective by the preceding proof. This, rather than the subsystem alone, proves the source's containment.

To recover the finite stages needed for the source's henselian conclusion, take any étale \(A\)-stage \(E\to C\). The proved henselian decomposition is \(E=E_1\times\cdots\times E_n\times B\), with the \(E_i\) finite local étale over \(A\) and \(\mathfrak m_A B=B\). A map to the nonzero local ring \(C\) sends exactly one of the finitely many product idempotents to 1 and the rest to 0. It cannot choose \(B\), because a finite expression \(1=\sum a_jb_j\), with \(a_j\in\mathfrak m_A\), would put 1 in \(\mathfrak m_C\). Thus it factors through a selected \(E_i\). The factor map to \(C\) is local: its contracted prime lies over \(\mathfrak m_A\), and a finite local algebra has only its maximal ideal over that ideal. For any two stages or any equality of elements, take a common stage in the original filtered presentation and then its selected finite factor. The transition maps factor through the previously selected factors because the product idempotents have the same 0 or 1 images in that local factor. This proves that the selected finite local stages are filtered and have colimit exactly \(C\), on both elements and equalities. The same henselian colimit proof now applies. The finite étale equivalence identifies these stages and their exact residue embeddings with the stated tower.

Finally keep the exact objects in 43540--43558. Write \(A=(R_{\mathfrak p})^h\), \(B=(S_{\mathfrak q})^h\), and assume the original residue map \(\kappa(\mathfrak p)\to\kappa(\mathfrak q)\) is an isomorphism. Use that given map to identify the residue fields, as the source does, and use its chosen \(\kappa^{sep}\). The source tensor product is \[
 T=B\otimes_A(R_{\mathfrak p})^{sh}
   =\mathop{\rm colim}_{\kappa'}\bigl(B\otimes_A A(\kappa')\bigr).
\] The forward comparison sends \(b\otimes[a]\) to \([b\otimes a]\); its inverse sends each stage tensor to the same tensor in \(T\). Balancing and the transition equations make these inverse maps on every generator.

Each original stage \(B\otimes_A A(\kappa')\) is finite étale over \(B\), with closed fibre exactly \(\kappa(\mathfrak q)\otimes_{\kappa(\mathfrak p)}\kappa'=\kappa'\) by the original residue isomorphism. The finite-algebra henselian decomposition over \(B\) therefore has exactly one nonzero local factor: every maximal ideal of a finite algebra over a local ring contracts to the maximal ideal, and the closed fibre here is a field. The stage is local and henselian. Its maximal ideal is \(\mathfrak m_B\) times the stage, and the transition maps are local, with the specified inclusions of residue fields. Hence \(T\) is a henselian local ind-étale \(B\)-algebra with residue \(\kappa^{sep}\); it is strictly henselian. It is also ind-étale over \(S_{\mathfrak q}\) by composition. The original tensor comparison \[
 T\longrightarrow(S_{\mathfrak q})^{sh},\qquad
 b\otimes a\longmapsto b\,f(a)
\] is local and induces the identity on the chosen residue closure. The unique local map in the reverse direction with the same residue identification exists by henselian lifting on its original ind-étale stages. Both composites are identities by uniqueness. This proves precisely the original formula, with its original residue-isomorphism hypothesis. No stronger residue-extension claim is recorded in this batch.

\subsubsection{Report dispositions and the deferred follow-on work}\label{algebra-proof-060-report-dispositions-and-the-deferred-follow-on-work}

The fifteen reports give nine groups and eight source operations:

\begin{enumerate}
\def\labelenumi{\arabic{enumi}.}
\tightlist
\item
  \texttt{OCC-00968}, \texttt{OCC-12274}, 43084--43089: reject the asserted defect. In ``Now that triple \ldots{} determines'', ``now'' is an adverb and ``that triple'' is its demonstrative subject. The text does not need a subordinate ``now that'' clause or a compulsory replacement of ``that'' by ``the''. The class of \(T\) proves the stated mathematical conclusion as above.
\item
  \texttt{OCC-00969}: add the missing terminal period at 43139.
\item
  \texttt{OCC-00970}, \texttt{OCC-12278}: at 43284 replace the undefined \(\mathfrak p'\) by the actual \(\mathfrak q'\). The original fibre and coequalizer maps above verify the referent.
\item
  \texttt{OCC-00971}, \texttt{OCC-12282}: at 43360 replace \(\mathfrak m_{S^h}\) by \(\mathfrak m_{S^{sh}}\), the maximal ideal of the actual codomain.
\item
  \texttt{OCC-00972}: join the given-diagram clause to its main clause at 43381--43383, inserting the comma and changing ``There'' to ``there''.
\item
  \texttt{OCC-00973}, \texttt{OCC-12283}: at 43422 refer to the category of triples. All embedding data remain as originally stated.
\item
  \texttt{OCC-00974}, \texttt{OCC-12284}: at 43441 remove the stray article and add the participial comma: ``Combining the maps above, we''.
\item
  \texttt{OCC-00975}, \texttt{OCC-12285}: at 43461 call the arrows morphisms of triples. They also are morphisms of the underlying pairs; this terminological edit does not assert that the original mathematical arrow was invalid.
\item
  \texttt{OCC-01463}: join the given-diagram clause at 43522--43524, inserting the comma and changing ``The'' to ``the''. The report's claim that there is a post-display period is inaccurate; the original has no such period. The accepted edit repairs the actual capitalization and clause joining.
\end{enumerate}

The earlier ind-étale lifting, maps-between-algebras and henselian results have been checked on these receiving source constructions. Their original source hypotheses suffice. Full reconstructions stay in this note; only the eight minimal operations enter the proposed correction layer. The optional wider consequences and further external-source reading are queued after completion of the Algebra review. No formal proof-assistant certification, chapter admission, TeX build, publication or novelty claim is made by this note.
